
\documentclass[12pt]{article}
%%%%%%%%%%%%%%%%%%%%%%%%%%%%%%%%%%%%%%%%%%%%%%%%%%%%%%%%%%%%%%%%%%%%%%%%%%%%%%%%%%%%%%%%%%%%%%%%%%%%%%%%%%%%%%%%%%%%%%%%%%%%%%%%%%%%%%%%%%%%%%%%%%%%%%%%%%%%%%%%%%%%%%%%%%%%%%%%%%%%%%%%%%%%%%%%%%%%%%%%%%%%%%%%%%%%%%%%%%%%%%%%%%%%%%%%%%%%%%%%%%%%%%%%%%%%
\usepackage{amsmath}
\usepackage[compat2]{geometry}

\setcounter{MaxMatrixCols}{10}
%TCIDATA{TCIstyle=LaTeX article (bright).cst}

%TCIDATA{OutputFilter=LATEX.DLL}
%TCIDATA{Version=5.50.0.2953}
%TCIDATA{<META NAME="SaveForMode" CONTENT="1">}
%TCIDATA{BibliographyScheme=Manual}
%TCIDATA{Created=Monday, May 23, 2011 14:12:05}
%TCIDATA{LastRevised=Monday, July 14, 2025 11:31:14}
%TCIDATA{<META NAME="GraphicsSave" CONTENT="32">}
%TCIDATA{<META NAME="DocumentShell" CONTENT="Exams and Syllabi\SW\Assignment">}

\setlength{\topmargin}{-1.0in}
\setlength{\textheight}{9.25in}
\setlength{\oddsidemargin}{0.0in}
\setlength{\evensidemargin}{0.0in}
\setlength{\textwidth}{6.5in}
\def\labelenumi{\arabic{enumi}.}
\def\theenumi{\arabic{enumi}}
\def\labelenumii{(\alph{enumii})}
\def\theenumii{\alph{enumii}}
\def\p@enumii{\theenumi.}
\def\labelenumiii{\arabic{enumiii}.}
\def\theenumiii{\arabic{enumiii}}
\def\p@enumiii{(\theenumi)(\theenumii)}
\def\labelenumiv{\arabic{enumiv}.}
\def\theenumiv{\arabic{enumiv}}
\def\p@enumiv{\p@enumiii.\theenumiii}
\pagestyle{plain}
\setcounter{secnumdepth}{0}
\parindent=0pt
\input{tcilatex}
\begin{document}


\begin{center}
\begin{equation*}
\FRAME{itbpF}{1.3214in}{0.9496in}{0in}{}{}{Figure}{\special{language
"Scientific Word";type "GRAPHIC";maintain-aspect-ratio TRUE;display
"USEDEF";valid_file "T";width 1.3214in;height 0.9496in;depth
0in;original-width 1.3552in;original-height 0.9669in;cropleft "0";croptop
"1";cropright "1";cropbottom "0";tempfilename
'R3IG7W00.bmp';tempfile-properties "XPR";}}
\end{equation*}%
\textbf{Electromagnetismo}

\emph{Gu\'{\i}a N}$%
%TCIMACRO{\U{b0}}%
%BeginExpansion
{{}^\circ}%
%EndExpansion
9:$ \emph{Fuentes de campo magn\'{e}tico}

\textit{Profesor: Arturo\ Fern\'{a}ndez P\'{e}rez}
\end{center}

\begin{enumerate}
\item (El cable coaxial) Un conductor s\'{o}lido de radio $a$ est\'{a}
sostenido por discos aislantes sobre el eje de un tubo conductor de radio
interior $b$ y radio exterior $c$, como muestra la Fig. 1. El conductor
central y el tubo transportan corrientes $I$ en sentidos opuestos. Determine
el m\'{o}dulo del campo magn\'{e}tico en:

\begin{description}
\item[a.] La regi\'{o}n $a<r<b$. \textbf{R: }$B=\frac{\mu _{0}I}{2\pi r}$

\item[b.] La regi\'{o}n $r>c$. \textbf{R: }$B=0$%
\begin{equation*}
\FRAME{itbpFU}{1.7469in}{2.0738in}{0in}{\Qcb{Figura 1}}{}{Figure}{\special%
{language "Scientific Word";type "GRAPHIC";maintain-aspect-ratio
TRUE;display "USEDEF";valid_file "T";width 1.7469in;height 2.0738in;depth
0in;original-width 3.6633in;original-height 4.3578in;cropleft "0";croptop
"1";cropright "1";cropbottom "0";tempfilename
'R3IG7W01.wmf';tempfile-properties "XPR";}}
\end{equation*}
\end{description}

\item Tres alambres transportan una corriente $I=4$ [A] y se encuentran a
una distancia $d=15$ [cm] entre ellos, como se muestra en la Fig. 2. Calcule
la fuerza magn\'{e}tica por unidad de longitud sobre cada alambre. \textbf{%
R: }$\frac{F_{1}}{l}=1,07\times 10^{-5}$ [N/m] hacia arriba, $\frac{F_{3}}{l}%
=1,07\times 10^{-5}$ [N/m] hacia abajo, $\frac{F_{2}}{l}=0$ [N]. \textbf{\ }%
\begin{equation*}
\FRAME{itbpFU}{2.2502in}{1.3154in}{0in}{\Qcb{Figura 2}}{}{Figure}{\special%
{language "Scientific Word";type "GRAPHIC";maintain-aspect-ratio
TRUE;display "USEDEF";valid_file "T";width 2.2502in;height 1.3154in;depth
0in;original-width 4.248in;original-height 2.4708in;cropleft "0";croptop
"1";cropright "1";cropbottom "0";tempfilename
'R3IG7X02.wmf';tempfile-properties "XPR";}}
\end{equation*}

\item Un alambre recto y largo $AB$, como se muestra en la Fig. 3, conduce
una corriente de $14$ [A]. La espira rectangular cuyos lados largos son
paralelos al alambre conducen una corriente de $5$ [A]. Determine el momento
magn\'{e}tico de la espira y la fuerza que el campo magn\'{e}tico del
alambre ejerce sobre cada lado de la espira. \textbf{R: }Lado superior $%
F=1,08\times 10^{-4}$ [N] que lo atrae hacia el alambre recto. Lado inferior 
$F=2,8\times 10^{-5}$ [N] que lo repele del alambre recto. Lado izquierdo $%
F=1,89\times 10^{-5}$ [N] dirigida hacia la izquierda. Lado derecho $%
F=1,89\times 10^{-5}$ [N] dirigida hacia la derecha. 
\begin{equation*}
\FRAME{itbpFU}{2.6048in}{1.9484in}{0in}{\Qcb{Figura 3}}{}{Figure}{\special%
{language "Scientific Word";type "GRAPHIC";maintain-aspect-ratio
TRUE;display "USEDEF";valid_file "T";width 2.6048in;height 1.9484in;depth
0in;original-width 4.3734in;original-height 3.2629in;cropleft "0";croptop
"1";cropright "1";cropbottom "0";tempfilename
'R3IG7X03.wmf';tempfile-properties "XPR";}}
\end{equation*}

\item En la Fig. 5, el cable que conduce una corriente $I=5$ $[A]$ forma un
arco circular de radio $R=4$ [cm]$.$ Determine el campo magn\'{e}tico en el
origen. \textbf{R: }$B=1,96\times 10^{-5}$ [T] entrando al plano de la p\'{a}%
gina.%
\begin{equation*}
\FRAME{itbpFU}{2.7466in}{2.1646in}{0in}{\Qcb{Figura 5}}{}{Figure}{\special%
{language "Scientific Word";type "GRAPHIC";maintain-aspect-ratio
TRUE;display "USEDEF";valid_file "T";width 2.7466in;height 2.1646in;depth
0in;original-width 3.8311in;original-height 3.013in;cropleft "0";croptop
"1";cropright "1";cropbottom "0";tempfilename
'R3IG7X05.wmf';tempfile-properties "XPR";}}
\end{equation*}%
\pagebreak

\item Dos conductores muy extensos conducen corrientes que se detallan en la
Fig. 6. Determine en qu\'{e} puntos el campo magn\'{e}tico es cero. \textbf{%
R:} En todos los puntos de la recta $y=\frac{3}{10}x$%
\begin{equation*}
\FRAME{itbpFU}{1.8948in}{1.7288in}{0in}{\Qcb{Figura 6}}{}{Figure}{\special%
{language "Scientific Word";type "GRAPHIC";maintain-aspect-ratio
TRUE;display "USEDEF";valid_file "T";width 1.8948in;height 1.7288in;depth
0in;original-width 3.1202in;original-height 2.8461in;cropleft "0";croptop
"1";cropright "1";cropbottom "0";tempfilename
'R3IG7X06.wmf';tempfile-properties "XPR";}}
\end{equation*}

\item El cable de la Fig. 7 conduce una corriente $I=3$ [A] y tiene una secci%
\'{o}n curva de radio $R=7$ [cm]. Determine el campo magn\'{e}tico en el
punto P. \textbf{R: }$B=-1,35\times 10^{-5}$ [T] entrando al plano de la p%
\'{a}gina. 
\begin{equation*}
\FRAME{itbpFU}{2.2943in}{1.1813in}{0in}{\Qcb{Figura 7}}{}{Figure}{\special%
{language "Scientific Word";type "GRAPHIC";maintain-aspect-ratio
TRUE;display "USEDEF";valid_file "T";width 2.2943in;height 1.1813in;depth
0in;original-width 4.5403in;original-height 2.3246in;cropleft "0";croptop
"1";cropright "1";cropbottom "0";tempfilename
'R3IG7X07.wmf';tempfile-properties "XPR";}}
\end{equation*}

\item Una espira circular de radio $R=10$ [cm] conduce una corriente $%
I_{2}=1 $ [A] en sentido horario, como muestra la Fig. 8. El centro de la
espira se encuentra a una distancia $D=20$ [cm] de un cable recto muy
extenso. Determine la magnitud y direcci\'{o}n de la corriente $I_{1}$ si el
campo magn\'{e}tico en el centro de la espira es cero. \textbf{R: }$I=6,28$
[A] circulando de izquierda a derecha.%
\begin{equation*}
\FRAME{itbpFU}{2.4414in}{1.5549in}{0in}{\Qcb{Figura 8}}{}{Figure}{\special%
{language "Scientific Word";type "GRAPHIC";maintain-aspect-ratio
TRUE;display "USEDEF";valid_file "T";width 2.4414in;height 1.5549in;depth
0in;original-width 4.6648in;original-height 2.9611in;cropleft "0";croptop
"1";cropright "1";cropbottom "0";tempfilename
'R3IG7X08.wmf';tempfile-properties "XPR";}}
\end{equation*}

\item Un electroim\'{a}n se compone de un solenoide que tiene 60 vueltas de
cable por cm, un radio $R=3$ [cm] y longitud $l=80$ $[cm],$ enrollado sobre
una aleaci\'{o}n de silicio y acero que posee permeabilidad magn\'{e}tica
relativa $K_{m}=\frac{\mu _{m}}{\mu _{0}}=5200$, como muestra la Fig. 9. \
Si por el solenoide circula una corriente de $3$ [A] en sentido antihorario,
determine $\vec{H}$, $\vec{M}$ y $\vec{B}$ en el centro del solenoide y el
momento magn\'{e}tico del solenoide. \textbf{R: }Todos los vectores son
paralelos y apuntan hacia la derecha de la p\'{a}gina. $H=18000$ [A/m], $%
M=9,36\times 10^{7}\,$[A/m], $B=117,64$ [T], $\mu =40,72$ [Am$^{2}$]%
\begin{equation*}
\FRAME{itbpFU}{2.5356in}{1.5126in}{0in}{\Qcb{Figura 9}}{}{Figure}{\special%
{language "Scientific Word";type "GRAPHIC";maintain-aspect-ratio
TRUE;display "USEDEF";valid_file "T";width 2.5356in;height 1.5126in;depth
0in;original-width 4.3734in;original-height 2.5962in;cropleft "0";croptop
"1";cropright "1";cropbottom "0";tempfilename
'R3IG7Y09.wmf';tempfile-properties "XPR";}}
\end{equation*}
\end{enumerate}

\end{document}
